\documentclass[10pt,a4paper]{amsart}
\usepackage[margin=19mm]{geometry}
\usepackage{amsmath,amssymb,mathtools}
\usepackage[scr=rsfs,cal=euler]{mathalfa}
\usepackage{xfrac}
\usepackage[unicode,hidelinks,bookmarks=false]{hyperref}
\newcommand{\ie}{i.e.\ }
\newcommand{\cf}{cf.\ }
\newcommand{\resp}{respectively\ }
\newcommand{\Subsection}{\subsection}
\providecommand{\nobreakdash}{\nobreak\hbox{-}}
\setlength{\emergencystretch}{2em}
\multlinegap0pt
\usepackage{xcolor}
\usepackage{amssymb}
\usepackage{mathtools}
\usepackage[shortlabels]{enumitem}
\usepackage{float}
\usepackage{tikz}
\newtheorem{thm}{Theorem}
\newtheorem{lem}{Lemma}
\title{Two-Generator Discrete and Faithful Subgroups of Tree Automorphisms}
\author{Yukun Du}
\date{Source: arXiv:2606.06824v2 (CC BY 4.0)}
\begin{document}
\maketitle

\noindent\emph{Source and licence.} Two-Generator Discrete and Faithful Subgroups of Tree Automorphisms. Version arXiv:2606.06824v2 (CC BY 4.0), distributed by arXiv under the Creative Commons Attribution 4.0 International licence, http://creativecommons.org/licenses/by/4.0/, as recorded in the arXiv licence metadata for this exact version and shown on the abstract page https://arxiv.org/abs/2606.06824v2. Full licence notice: \texttt{licenses/arxiv-2606.06824-CC-BY-4.0.txt}.\par\smallskip

\noindent\textit{Editorial scope.} Counted unit: the article's thm thm:HNN\_rigidity (source0.tex lines 449--451). The article's complete original proof of it is delivered with it (source0.tex lines 514--623, one \texttt{proof} environment of 110 lines, which the article places away from the statement). No mathematics is written, repaired or re-derived: the statement and the proof are the article's own lines, in the article's own notation, and no retained line is substituted or re-flowed.  Retained uncounted as the source-local context the proof needs: lines 453--455; lines 486--496.  Nothing was omitted from any retained span, so the package carries no omission record.  No retained line contains a citation, so the wrapper carries no bibliography and no bibliography entry.  No retained line contains a figure environment, an image-inclusion command or drawing code, so no image is delivered and the package declares no asset dependency.  Editorial formatting only, in addition: an AMS \texttt{amsart} preamble that restates the article's own declarations and macros used by the retained lines, namely the environments \texttt{thm}, \texttt{lem} on separate counters, together with the standard packages those lines need.  The retained environments print in this document as Theorem 1, Lemma 1 in order of appearance, whereas the article prints the same items with the numbering of its own full document.  The complete native source of the article as distributed in the arXiv e-print for this exact version is bundled unchanged as \texttt{sources/202216/source0.tex}.
\begin{lem}\label{lem:pres}
    The group $G$ satisfying the requirements in Theorem \ref{thm:HNN_rigidity} is a $2$-group. In addition, if $\lvert G\rvert = 2^n$, then $G$ is generated by $n$ involutions $t_1,\dots, t_n$, such that the HNN isomorphism satisfies $\varphi(t_i) = t_{i+1}$.
\end{lem}

\begin{lem}\label{lem:comm}
    Each element in $G$ is uniquely expressed as a product,
    \[
    g = t_{i_1}\dots t_{i_k},\ 1\leq i_1<\dots < i_k\leq n.
    \]
    Moreover, the group $G$ admits commutator relations,
    \[
    [t_i,t_{i+l}] = t_{i+1}^{\epsilon(l,1)}t_{i+2}^{\epsilon(l,2)}\dots t_{i+l-1}^{\epsilon(l,l-1)},
    \]
    for $\epsilon(l,1)$,\dots, $\epsilon(l,l-1)\in \{0,1\}$.
\end{lem}

\begin{thm}\label{thm:HNN_rigidity}
    Suppose $\Gamma = G*_{\varphi}$ is a faithful $2$-HNN extension of a finite group $G$—that is, $H_1,H_2 < G$ with $\lvert G:H_1\rvert = \lvert G:H_2\rvert = 2$, $\varphi: H_1 \to H_2$ is an isomorphism, and no non-trivial subgroup of $H_1 \cap H_2$ is preserved by $\varphi$. Then $\Gamma$ is isomorphic to an HNN extension constructed as above. In particular, every element in $G$ has order at most $4$.
\end{thm}

\begin{proof}[Proof of Theorem \ref{thm:HNN_rigidity}]
    We will prove by induction that the relators of such a group satisfy the presentation in Theorem \ref{thm:HNN_rigidity}. The base case is clear: for orders $2^1$ and $2^2$, $G = C_2$ and $G = V_4$ are the only eligible groups.
    
    For convenience, we assume that the theorem holds for groups of order $2^{n-1}$, and we prove the case $\lvert G\rvert = 2^n$ under this assumption. By Lemma \ref{lem:pres}, the subgroups $H_1 = \langle t_1,\dots, t_{n-1}\rangle$ and $H_2 = \langle t_2,\dots, t_n\rangle < G$ satisfy the assumption in the theorem. Thus, we can assume that there exists $0\leq k_0\leq \frac{n-1}{3}$ such that $t_i$ and $t_{i+j}$ commute for any $j\leq n-2-k_0$, and for $n-1-k_0\leq l\leq n-2$,
    \[
    [t_i,t_{i+l}] = t_{i+k_0}^{\epsilon(l,k_0)}\dots t_{i+l-k_0}^{\epsilon(l,l-k_0)},
    \]
    where $\sum_j\epsilon(n-1-k_0,j)>0$. Denote
    \[
    k_+ = \max_l(\min_{\epsilon(l,k)=1}k),\ k_- = \max_l(l - \max_{\epsilon(l,k) = 1} k),
    \]
    the induction hypothesis implies that $k_+,k_-\geq k_0$, while the case $\lvert G\rvert = 2^n$ further requires that $k_+,k_-\geq k_0+1$.
    
    Suppose that $k_+ = k_0$; then there exists some $l_0\geq n-1-k_0$ such that
    \[
    (t_1t_{l_0+1})^2 = t_{k_1'}\dots t_{k_j'},
    \]
    with $k_0+1=k_1'\leq\dots\leq k_j'\leq l-k_0+1$.
    
    Now, since $H_1$ is a normal subgroup of $G$, we have an automorphism:
    \[
    \psi:H_1\to H_1,\ t\mapsto t_ntt_n.
    \]
    The induction hypothesis implies that $\psi(t_j) = t_j$ for $j\geq k_0+2$, and
    \[
    \psi(t_j) = t_j t_{j+k_0}^{\epsilon(n-j,k_0)}\dots t_{n-k_0}^{\epsilon(n-j,n-j-k_0)},
    \]
    for $j=2,\dots, k_0+1$. In particular, the fact that $\sum_l\epsilon(n-1-k_0,l)>0$ in the definition of $k_0$ implies that $[t_{k_0+1},t_n]$ is non-trivial, or
    \[
    \psi(t_{k_0+1}) = t_{k_0+1}t_{k_1}\dots t_{k_i} \neq t_{k_0+1},
    \]
    with $2k_0+1\leq k_1<\dots< k_i\leq n-k_0$ and $i\geq 1$.
    
    Combine the facts above:
    \[
    \begin{split}
        & (\psi(t_1)t_{l_0+1})^2 = (\psi(t_1)\psi(t_{l_0+1}))^2 = \psi((t_1t_{l_0+1})^2) \\
        = & \psi(t_{k_0+1} t_{k_2'}\dots t_{k_j'}) = \psi(t_{k_0+1})t_{k_2'}\dots t_{k_j'} \\
        \neq & t_{k_0+1}t_{k_2'}\dots t_{k_j'} = (t_1t_{l_0+1})^2.
    \end{split}
    \]
    While the induction assumption does not allow $\psi(t_1)$ to fit into the pattern, Lemma \ref{lem:comm} guarantees that $t_1\psi(t_1)$ can be expressed as a word in the letters $t_2,\dots,t_{n-1}$ in ascending order. If $t_1\psi(t_1)$ is a word involving only the letters $t_{k_0+1},\dots,t_{n-1}$, then the commutativity of $t_{l_0+1}$ with $t_1\psi(t_1)$ implies that $(\psi(t_1)t_{l_0+1})^2 = (t_1t_{l_0+1})^2$, a contradiction. Therefore, if we let
    \[
        m = \min_{\epsilon(n-1,k)=1}k,
    \]
    then $m\leq k_0-1$, that is, $m+n-k_0\leq n-1$. The commutator $[t_1,t_{m+n-k_0}]$ can be expressed as a word in the letters $t_{k_0+1},\dots,t_{n-k_0-1}$.

    If this word expression of $[t_1,t_{m+n-k_0}]$ does not contain $t_{k_0+1}$, then it is fixed by the automorphism $\psi$. The element $t_{m+n-k_0}$ is also in this fixed-point subgroup:
    \[
        [t_1,t_{m+n-k_0}] = \psi([t_1,t_{m+n-k_0}]) = [\psi(t_1),\psi(t_{m+n-k_0})] = [\psi(t_1),t_{m+n-k_0}].
    \]
    By the assumption, $\psi(t_1) = t_1t_{m+1}\prod_{j\geq m+1}t_{j+1}^{\epsilon(n-1, j)}$, and $t_{m+n-k_0}$ commutes with every such $t_{j+1}$ as $j+1\geq m+2$. As a result,
    \[
        [t_{m+1},t_{m+n-k_0}] = 1,
    \]
    contradicting the fact that $\sum_l\epsilon(n-1-k_0,l)>0$.
    
    For the case that $t_{k_0+1}$ is contained in the commutator $[t_1,t_{m+n-k_0}]$, we need a further lemma:
    \begin{lem}\label{lem:not_contain}
        In the setting above, we have $\epsilon(n-k_0-1, n-2k_0-1) = 0$. That is, $t_{n-k_0}$ does not appear as a letter in the word expression of $\psi(t_{k_0+1})$. 
    \end{lem}
    \begin{proof}
        Since $t_1$ commutes with $t_{k_0+1}$, we have
        \[
            [\psi(t_1),\psi(t_{k_0+1})] = \psi([t_1,t_{k_0+1}]) = 1.
        \]
        If $t_{n-k_0}$ appears in the word expression of $\psi(t_{k_0+1})$, then any other factor lies between $t_{k_0+1}$ and $t_{n-k_0-1}$, hence lies in the center of $H_1 = \langle t_1,\dots, t_{n-1}\rangle$ and therefore commutes with $\psi(t_1)$. Therefore,
        \[
        [\psi(t_1),t_{n-k_0}] = 1.
        \]
        However, every factor in $\psi(t_1)$ except for $t_1$ itself commutes with $t_{n-k_0}$, implying $[t_1,t_{n-k_0}] = 1$, which contradicts the definition of $k_0$.
    \end{proof}
    Returning to the proof of Theorem \ref{thm:HNN_rigidity}. Since $t_{k_0+1}$ plays a role in $[t_1,t_{m+n-k_0}]$, a slightly different relation holds:
    \[
        [\psi(t_1),t_{m+n-k_0}] = \psi([t_1,t_{m+n-k_0}]) = t_{k_0+1}\psi(t_{k_0+1})[t_1,t_{m+n-k_0}].
    \]
    Similar to the previous case, the fact that $t_{m+n-k_0}$ commutes with every factor in $\psi(t_1)$ except for $t_1$ and $t_{m+1}$ implies
    \[
        [t_{m+1},t_{m+n-k_0}] = t_1t_{k_0+1}\psi(t_{k_0+1})t_1.
    \]
    However, Lemma \ref{lem:not_contain} implies that the index of each letter in the word expression of $\psi(t_{k_0+1})$ is strictly less than $n-k_0$, so $\psi(t_{k_0+1})$ commutes with $t_1$. Hence,
    \[
    [t_{m+1},t_{m+n-k_0}] = t_{k_0+1}\psi(t_{k_0+1}) = [t_{k_0+1},t_n],
    \]
    or
    \[
    \prod_{j=k_0}^{n-2k_0-1}t_{m+1+j}^{\epsilon(n-k_0-1,j)} = \prod_{j=k_0}^{n-2k_0-1}t_{k_0+1+j}^{\epsilon(n-k_0-1,j)}
    \]
    which contradicts the definition of $k_0$, the fact that $m<k_0$, and the unique ascending-letter expression of elements in $G$. This shows that $k_+\geq k_0+1$. A similar argument for the $t_1$-conjugation automorphism on $H_2$ implies that $k_-\geq k_0+1$.
    
    It remains to show that the commutator $[t_1,t_n] = (t_1t_n)^2$ also fits the pattern of the other commutators; that is, it is a word in ascending letters among $t_{k_0+2}$ through $t_{n-k_0-1}$. As mentioned earlier, a rough characterization by Lemma \ref{lem:comm} implies that $[t_1,t_n]$ is a word in letters from $t_2$ to $t_{n-1}$.

    First, suppose the word expression of $[t_1,t_n]$ contains letters from $t_2$ to $t_{k_0}$. In the definition of $m = \min_{\epsilon(n-1,k) = 1} k$ above, this implies $m\leq k_0 - 1$. The fact that $k_+\geq k_0+1$ implies that the word expression of $[t_1,t_{m+n-k_0}]$ does not contain the letter $t_{k_0+1}$, leading to the same contradiction that $[t_{m+1},t_{m+n-k_0}] = 1$. Therefore, $[t_1,t_n]$ does not contain letters from $t_2$ to $t_{k_0}$, and similarly, it does not contain letters from $t_{n-k_0+1}$ to $t_{n-1}$.
    
    Now suppose $[t_1,t_n]$ contains the letter $t_{k_0+1}$, that is,
    \[
        \psi(t_1) = t_1t_{k_0+1}w,
    \]
    for some word $w$ in letters with indices $\geq k_0+2$ and $<n$. Since conjugation by $t_n$, namely $\psi$, has order $2$, we have
    \[
        t_1 = \psi^2(t_1) = \psi(t_1)\psi(t_{k_0+1})\psi(w) = t_1t_{k_0+1}wt_{k_0+1}w'w,
    \]
    where $w' = t_{k_0+1}\psi(t_{k_0+1})$ is non-identity by the definition of $k_0$. Now, the fact that the word $w$ has indices between $k_0+2$ and $n-1$ implies that $w^2 = 1$ and $[w,t_{k_0+1}] = 1$, so
    \[
        1 = t_{k_0+1}wt_{k_0+1}w'w = ww'w\implies w' = w^2 = 1,
    \]
    a contradiction. Similarly, $t_{n-k_0}$ is not in the word expression of $[t_1,t_n]$, proving our claim.
    
    It is now straightforward that the square of every element in $G$ is contained in the Abelian subgroup $\langle t_{k+1},\dots, t_{n-k}\rangle$, and the fourth power equals the identity.
\end{proof}

\end{document}
